\scanpage{49}

% 1–4
Soit $\mathbb P_{\bullet}$ la catégorie des complexes de $k$-chaînes $L_{\bullet}$ qui sont libres de r. f. en chaque dim, et qui n'ont qu'un nb fini de composantes non nulles. Une suite

% 1–7, left margin
% slanted
\margin{Il vaut mieux construire $\mathbb P$ ``en l'air'', puis le ``réaliser'' \deleted{\unsure{quatre}} $\mathbb P_{\bullet},\mathbb P_{*},\mathbb P^{\bullet},\mathbb P^{*}$ comme pour $\mathbb K$.}

% 5
\[
 0\longrightarrow L'_{\bullet}\longrightarrow L_{\bullet}\longrightarrow L''_{\bullet}\longrightarrow 0
\]

% 6–8
\inserted{d'hom dans $\mathbb P_{\bullet}$} est dite exacte spéciale, si $\forall n\geq0$, la suite $0\to L'_n\longrightarrow L_n\longrightarrow L''_n\longrightarrow0$ est exacte. Soit $\Sigma_{\bullet}$ la classe de ces suites.

% 8–11
Ainsi $\mathbb P_{\bullet}$ est muni de $\mathbb E_{\bullet}^{\bullet}$ (complexe de cochaînes) et de $\Sigma_{\bullet}$. On va voir que pour ces données, il est sol. d'un pb universel.

% 11, left margin
\margin{!}

% 12–13
\deleted{Soit \unsure{$A$}} $B$ une cat. additive. Un ens. $\Sigma$ de suites

% 14, (x)
\[
 0\longrightarrow A'\xrightarrow{q} A\xrightarrow{p}A''\longrightarrow0 \tag*{(x)}
\]

% 15–16
dans $B$ est dit ``une classe de suites exactes'' si elle satisfait aux conditions

% 17–20
Sex 1) Si $(x)\in\Sigma$, alors $A'\overset{\sim}{\longrightarrow}\operatorname{Ker}p$, $\operatorname{Coker}q\overset{\sim}{\longrightarrow}A''$, i.e. les suites $0\longrightarrow A'\longrightarrow A\longrightarrow A''$ et $\deleted{\unsure{0\to}}A'\longrightarrow A\longrightarrow A''\longrightarrow0$ sont exactes (i.e. $\forall B\in\operatorname{Ob}B$, les suites de groupes ab.

% 21–22
\[
 \left\{\begin{aligned}
 0&\longrightarrow\operatorname{Hom}(B,A')\longrightarrow\operatorname{Hom}(B,A)\longrightarrow\operatorname{Hom}(B,A'')\\
 \deleted{\text{et}}\quad0&\longrightarrow\operatorname{Hom}(A'',B)\longrightarrow\operatorname{Hom}(A,B)\longrightarrow\operatorname{Hom}(A',B)
 \end{aligned}\right.
\]

% 23
sont exactes

% 24–25
Sex 2) $\Sigma$ stable par iso, par sommes finies et contient les suites de la forme

% 26–27
\[
 \begin{aligned}0&\longrightarrow A\xrightarrow{\mathrm{id}} A\longrightarrow0\longrightarrow0\\
 0&\longrightarrow0\longrightarrow A\xrightarrow{\mathrm{id}} A\longrightarrow0\end{aligned}
\]

% 28
(donc aussi les suites de la forme

% 29
\[
 0\longrightarrow A\xrightarrow{\mathrm{inj}}A\times B\xrightarrow{\mathrm{pr.}}B\longrightarrow0\qquad)
\]

% archive, bottom left
% archive sheet number: 48
